\documentclass[11pt,a4paper]{extarticle}

\usepackage{kotex}
\usepackage[margin=18mm]{geometry}
\usepackage{amsmath,amssymb,mathtools}
\usepackage{array,booktabs,longtable,tabularx}
\usepackage{enumitem}
\usepackage[hidelinks]{hyperref}
\usepackage{xurl}

\setlength{\parindent}{0pt}
\setlength{\parskip}{5pt}
\setlist[itemize]{leftmargin=18pt,itemsep=2pt,topsep=2pt}
\setlist[enumerate]{leftmargin=22pt,itemsep=3pt,topsep=3pt}
\setcounter{tocdepth}{2}

\newcommand{\Hash}{\mathrm{H}}
\newcommand{\Poseidon}{\mathrm{Poseidon2}}
\newcommand{\HToF}{\mathrm{HashToField}}
\newcommand{\CM}{\mathrm{CM}}
\newcommand{\DocumentHash}{\mathrm{DocumentHash}}
\newcommand{\List}[1]{\mathrm{List}_{#1}}
\newcommand{\Map}[1]{\mathrm{Map}_{#1}}

\title{zkDPP Development Specification}
\author{Canonical Implementation Specification}
\date{2026-07-16}
\renewcommand{\contentsname}{목차}

\begin{document}
\maketitle

\section*{문서 지위}

이 문서는 \texttt{Protocol Spec.tex}의 protocol semantics를
Go, gnark, Solidity와 Foundry로 재현하기 위한 구현 명세다. 객체의 의미, Relation과
state transition은 Protocol Specification이 우선하며, encoding, artifact, ABI,
storage와 test의 구체값은 이 문서가 우선한다. 미확정 항목은 구현자가 임의로 결정하지
않고 Open Decisions에 기록한다. 결정 이유는
\texttt{Decision Log.tex}에서 관리한다.

\tableofcontents
\newpage

\section{확정 개발 기준}

\subsection{기술 Stack}

\begin{center}
\begin{tabularx}{\textwidth}{@{}lX@{}}
\toprule
구분 & 확정 선택 \\
\midrule
Client 및 prover & Go \\
Circuit framework & gnark v0.15.0 \\
Crypto dependency & gnark-crypto v0.20.1 \\
Proof system & PLONK \\
Curve & BLS12-381 \\
Circuit field & BLS12-381 scalar field $\mathbb{F}_r$ \\
Polynomial commitment & KZG \\
Circuit-friendly hash & gnark BLS12-381 default Poseidon2 \\
Smart Contract & Solidity \\
Contract test & Foundry v1.7.1, Prague EVM \\
POC KZG setup & gnark \texttt{test/unsafekzg.NewSRS} \\
\bottomrule
\end{tabularx}
\end{center}

\texttt{unsafekzg}는 local POC에만 사용한다. Production SRS 또는 trusted setup
artifact로 취급하지 않는다.

\subsection{Circuit 집합}

\[
\mathcal{I}=\{
\mathrm{Entry},\mathrm{Transfer},\mathrm{Proceed},\mathrm{Recall},
\mathrm{Merge},\mathrm{Split},\mathrm{Process},\mathrm{Exit}
\}
\]

각 $i\in\mathcal{I}$에 대해 별도의 constraint system, proving key와 verifying key를
생성한다. Relation의 정확한 statement와 witness는 Protocol Specification의 동일한
Event 절을 따른다.

\subsection{현재 POC 경계}

\begin{itemize}
  \item 저장소 root의 \path{poc/}와 \path{contracts/}는 Entry와 고정 3-to-3 Process를
        검증한 historical baseline이다. 최신 Protocol 구현을 위해 이 두 폴더를 수정하지 않는다.
  \item 새 구현의 source는 \path{poc-v2/}이며 Protocol Specification의 8개 Event를
        대상으로 한다. 기능 확인용 기본 profile은 $\ell=3$, $d=32$,
        $M_{\max}=N_{\max}=3$이고 Process는 $1\leq m,n\leq3$ active prefix와 inactive
        zero padding을 구현한다.
  \item Merge는 fixed 2-to-1, Split은 fixed 1-to-2 circuit이다.
  \item 기능 확인용 기본 Process는 $M_{\max}=N_{\max}=3$이다. 성능 측정은
        $M_{\max}=N_{\max}=a$, $a\in\{1,2,3\}$인 별도 circuit을 compile하고
        1-to-1, 2-to-2와 3-to-3만 비교한다. 비대칭 arity는 측정하지 않는다.
  \item 기능 확인용 기본 StateVector 길이는 $\ell=3$이다. 성능 측정은
        $\ell\in\{0,1,2,3\}$인 별도 circuit을 compile한다.
  \item Pedersen commitment와 circuit 내부 Pedersen opening 검증은 사용하지 않는다.
  \item DocumentInfo와 DocumentHash의 일치는 circuit과 contract에서 검증하지 않는다.
  \item Smart Contract는 provenance edge event를 emit하지 않는다. 다만 Client가
        Merkle witness를 구성할 수 있도록 Tree append index event는 emit한다.
  \item Transfer는 $q=0$이어도 sender의 change commitment를 항상 생성한다.
  \item POC의 공급망 Event와 Recall resolution Event 함수는 caller를 제한하지 않는다.
        Transaction submitter는 note owner일 필요가 없고, note 또는 voucher 소비 권한은
        Circuit relation으로 검증한다. \texttt{RegisterUser}만 호출자의 Ethereum
        \texttt{msg.sender}를 편의용 $(addr,pk_{own})$ metadata에 연결한다.
\end{itemize}

Historical baseline의 재현 명령과 기존 측정값은 \path{poc/README.md}를 따른다.
새 POC의 구현 상태, 재현 명령과 측정값은 \path{poc-v2/README.md}에서 관리한다.
두 POC 사이에서 generated verifier, key, fixture 또는 benchmark 결과를 공유하지 않는다.

\subsection{기존 POC 재사용 경계}

기존 파일을 직접 import하여 두 POC를 결합하지 않는다. 검토를 통과한 공통 구현만
\path{poc-v2/}로 복사하고 새 module 안에서 독립적으로 유지한다.

\begin{longtable}{@{}p{0.22\textwidth}p{0.31\textwidth}p{0.39\textwidth}@{}}
\toprule
분류 & 대상 & 처리 원칙 \\
\midrule
\endhead
재사용 & gnark와 gnark-crypto version, Go version & 동일 version으로 고정하되 새
module의 dependency로 선언 \\
재사용 & Poseidon2 native/circuit hash와 compressor & Protocol reference vector를
다시 통과한 소스만 복사 \\
재사용 & Depth-32 native/circuit Merkle Tree & $MT$와 $rvMT$의 공통 component로
사용하되 두 Tree의 storage와 root set은 분리 \\
재사용 & Poseidon2 Solidity generator, artifact writer와 manifest 구조 & relation
registry와 public input manifest는 8개 Event 기준으로 다시 작성 \\
재사용 & Docker Foundry environment와 \texttt{foundry.toml} 기본값 & 새 POC 내부의
Contract와 test만 mount하고 실행하도록 경로 변경 \\
미재사용 & 기존 Entry/Process/SHA circuit과 scenario & 최신 Relation을 Protocol
Specification에서 다시 전사 \\
미재사용 & 기존 main Contract, generated verifier, key, fixture & 23-input Process와
단일 root 등 이전 statement에 binding되어 있으므로 폐기 \\
미재사용 & 기존 benchmark 수치와 SHA profile & historical baseline으로만 보존하고
새 POC 결과에 혼합하지 않음 \\
\bottomrule
\end{longtable}

\section{구현 Component}

\begin{longtable}{@{}p{0.22\textwidth}p{0.70\textwidth}@{}}
\toprule
Component & 책임 \\
\midrule
\endhead
Client & note와 secret 보관, witness 구성, proof 생성, transaction input 준비 \\
Circuit & commitment opening, ownership, nullifier, membership, range와 Event relation 검증 \\
Setup tool & circuit compile, POC SRS 생성, PLONK key 생성, verifier와 artifact export \\
Smart Contract & registry와 root 검사, proof verification, duplicate/spent 검사, tree와 mapping 갱신 \\
Off-chain 전달 & note opening, $o$, $o_{rv}$와 필요한 product data를 권한 있는 상대방에게 전달 \\
Test harness & Go proof 검증과 Solidity verifier의 positive/negative case 재현 \\
\bottomrule
\end{longtable}

Off-chain 전달 protocol의 암호화 방식은 현재 POC 범위에서 제외한다. 다만 receiver가
Proceed 또는 이후 Event를 실행하려면 필요한 opening을 실제로 전달받아야 한다.

\subsection{구현 준비도}

2장 이후의 구현 가능 상태는 다음과 같다. ``차단''은 해당 결정을 확정하기 전에는
canonical implementation을 고정하지 않는다는 뜻이다. 독립적인 package skeleton이나
최소 실험까지 금지한다는 뜻은 아니다.

\begin{longtable}{@{}p{0.18\textwidth}p{0.15\textwidth}p{0.59\textwidth}@{}}
\toprule
영역 & 상태 & 남은 결정 \\
\midrule
\endhead
Component & POC 가능 & Client, Circuit, Setup, Contract와 test harness의 책임이 확정됐다. 구체 package 이름은 구현 선택이다. \\
Data Encoding & POC 가능 & DocumentInfo는 NFC 정규화, UTF-8, 4-byte unsigned big-endian length prefix로 고정됐다. Poseidon2와 \texttt{fr.Hash}도 확정됐다. \\
Setup/Artifact & POC 가능 & Relation별 \texttt{unsafekzg.NewSRS(ccs)}와 artifact directory 및 manifest schema가 확정됐다. \\
Tree/Storage & POC 가능 & Depth 32, Azeroth식 raw leaf, append, accepted-root와 index event 정책이 확정됐다. Depth scaling은 수행하지 않는다. \\
Circuit & POC 가능 & 공통 hash/tree/public-input ordering과 Transfer/Split deterministic remainder가 확정됐다. Process는 fixed-direction public $P_{pub}$, public $A$, constant $D=10^9$과 $m,n$ 기반 internal selector를 사용한다. \\
Smart Contract & POC 가능 & POC caller policy, storage, path view 함수와 Relation별 verifier 배포 구조가 확정됐다. custom error 이름은 의미론을 바꾸지 않는 구현 선택이다. \\
Test & POC 가능 & cross-layer reference vector, Event별 positive/negative case와 benchmark 항목이 확정됐다. 구체 test vector 값은 구현 중 생성한다. \\
\bottomrule
\end{longtable}

\section{Data Encoding과 Layout}

\subsection{DocumentInfo와 DocumentHash}

POC의 logical object는 다음과 같다.

\[
\mathit{DocumentInfo}=(ProductName,LotId,Unit)
\]

\[
 b_x=\operatorname{UTF8}(\operatorname{NFC}(x))
\]

\[
\operatorname{EncodeText}(x)
=\operatorname{U32BE}(|b_x|)
\mathbin{\|}b_x
\]

\[
\begin{aligned}
\mathit{DocumentInfoBytes}
={}&\operatorname{EncodeText}(ProductName)\\
&\mathbin{\|}\operatorname{EncodeText}(LotId)\\
&\mathbin{\|}\operatorname{EncodeText}(Unit)
\end{aligned}
\]

\[
\DocumentHash=\mathrm{fr.Hash}(
\mathit{DocumentInfoBytes},
\varepsilon,1
)[0]
\]

\texttt{fr.Hash}는 gnark-crypto v0.20.1의 BLS12-381 scalar-field
Hash-to-Field API다. 내부적으로 SHA-256 기반 \texttt{expand\_message\_xmd}를
사용하지만, 단순한 \texttt{SHA256(msg) mod r} 계산은 아니다. Receiver는
off-chain에서 동일한 serialization과 hash를 수행해 DocumentHash를 확인한다.
두 번째 API 인자 DST에는 빈 byte string을 전달하며 별도 hash 용도 구분값은
사용하지 않는다. Layout version은 넣지 않는다. 각 문자열은 Unicode NFC로 정규화하고
유효한 UTF-8 bytes로 encode한다. UTF-8 rune 수가 아니라 byte 수를 unsigned 32-bit
big-endian으로 먼저 기록한다. UTF-8이 유효하지 않거나 byte 길이가 $2^{32}-1$을
초과하면 encoding을 거부한다.

Go 구현은 \path{golang.org/x/text/unicode/norm}의 NFC 정규화와
\path{encoding/binary.BigEndian.PutUint32}를 사용한다. Field 순서는
\texttt{ProductName}, \texttt{LotId}, \texttt{Unit}으로 고정한다. Empty string과
NUL byte는 length prefix로 경계가 보존되므로 encoding 자체를 허용한다. POC는
DocumentInfo field의 비어 있음이나 business 의미를 별도로 검증하지 않는다.

\subsection{Quantity와 StateVector}

\[
q\in\{0,\ldots,2^{64}-1\}
\]

\[
\mathbf{s}=(s_1,\ldots,s_{\ell}),
\qquad
(s_1,s_2,s_3)=(total_{kg},carbon_{kgCO2e},recycled_{kg})
\]

기능 확인용 기본 POC에서 $\ell=3$이다. 성능 측정 profile은 다음 prefix layout을 사용한다.

\[
\begin{array}{c|l}
\ell & \mathbf{s} \\
\hline
0 & () \\
1 & (total_{kg}) \\
2 & (total_{kg},carbon_{kgCO2e}) \\
3 & (total_{kg},carbon_{kgCO2e},recycled_{kg})
\end{array}
\]

$q$는 integer count이고, 각 $s_k$는 unsigned 64-bit
integer로 제한한 nano scale $10^9$ fixed-point 값이다. 각 Event의 곱셈과 aggregate가
$\mathbb{F}_r$에서 wraparound하지 않도록 intermediate bound를 circuit별로 분석해야 한다.

\subsection{Owner와 Address}

별도 설명이 없는 한 이 문서의 Field는 BLS12-381 scalar field
$\mathbb{F}_r$를 뜻한다. Base field $\mathbb{F}_q$는 elliptic-curve point 좌표와
pairing 구현에서 사용한다.

\[
sk_{own}\xleftarrow{\$}\mathbb{F}_r,
\qquad
pk_{own}=\Hash(sk_{own}),
\qquad
addr=\Hash(pk_{own})
\]

\[
usk=sk_{own},
\qquad
upk=(addr,pk_{own})
\]

Holder는 등록된 $addr$을 알고 사용한다. Note 안에는 $addr$을 저장하지 않는다.
Spend circuit은 witness로 받은 $sk_{own}$에서 owner binding을 검증한다.
$pk_{own}$은 EC public key point가 아니라 Poseidon2로 도출한 public ownership
identifier다. $sk_{own}$과 $o$는 $\mathbb{F}_r$에서 sample하고,
$pk_{own}$, $addr$, $cm$은 Poseidon2 출력이므로 모두 $\mathbb{F}_r$ 원소다.

\texttt{RegisterUser}는 누구나 호출할 수 있다. 입력은 $addr$와 $pkOwn$이다.
Contract는 호출자의 Ethereum EOA인 \texttt{msg.sender}를 $(addr,pkOwn)$에 저장한다. POC에서는
$addr$ 중복, 동일 EOA 재등록과 $addr=\Hash(pkOwn)$을 Contract에서 검사하지 않으며,
같은 EOA가 다시 등록하면 최신 값으로 덮어쓴다. 이 mapping은 wallet 조회를 위한
off-circuit metadata다. Supply-chain Event proof의 owner authorization은 이 mapping이나
\texttt{msg.sender}가 아니라 Circuit의 $sk_{own}$ relation으로 결정된다. Public mapping의
자동 getter를 사용하고 registration event는 emit하지 않는다.

\subsection{Commitment와 Note}

Go, circuit과 note reconstruction에서 다음 flat ordering을 동일하게 사용한다.

\[
cm=\Poseidon(
\DocumentHash,q,s_1,\ldots,s_{\ell},addr,o
)
\]

\[
\mathit{note}=(cm,\DocumentHash,q,\mathbf{s},o)
\]

POC의 모든 $\Hash$와 $\CM$은 gnark v0.15.0이 의존하는 gnark-crypto v0.20.1의
BLS12-381 default Poseidon2 hasher를 사용한다. 구체 instance는 width 2,
full rounds 6, partial rounds 50, S-box $x^5$와 Merkle--Damg\aa rd construction이다.
입력 Field Element는 각 식에 적힌 왼쪽에서 오른쪽 순서로 hasher에 전달한다. 별도
해시 용도 구분값은 사용하지 않는다. 해시 용도 구분값은 같은 hash 함수를 서로 다른
protocol 목적으로 사용할 때 목적을 구분하기 위해 입력에 추가하는 고정값을 뜻한다.
DocumentHash용 \texttt{fr.Hash}도 별도 구분값을 사용하지 않으며, API가 요구하는
DST에는 빈 byte string을 전달한다.

다음 SHA-to-Field profile은 EVM SHA-256 precompile의 on-chain gas와 generated
Poseidon2 Solidity 연산의 gas를 같은 protocol 의미론에서 비교하기 위해 구현한 historical
benchmark다. 비교 실험은 종료했으며 이후 POC 구현과 benchmark에는 Poseidon2만 사용한다.
SHA profile은
DocumentHash용 \texttt{fr.Hash}와 다른 함수다. 입력
$x_1,\ldots,x_t\in\mathbb{F}_r$에 대해 다음 SHA-to-Field를 사용한다.

\[
H_{\mathrm{SHA}}(x_1,\ldots,x_t)=
\operatorname{OS2IP}\!\left(
\operatorname{SHA256}(\operatorname{I2OSP}_{32}(x_1)\|\cdots\|\operatorname{I2OSP}_{32}(x_t))
\right)\bmod r
\]

여기서 $\operatorname{I2OSP}_{32}$는 Field Element의 canonical unsigned 32-byte
big-endian encoding이다. Native Go는 SHA-256 digest를 BLS12-381 scalar field로
환원하고, circuit은 같은 byte ordering으로 SHA-256을 계산한 뒤 Field arithmetic으로
동일한 환원을 수행한다. Solidity는 \texttt{sha256} 결과를 동일한 $r$로 나눈 나머지를
사용한다. 두 Field 입력의 $H_{\mathrm{SHA}}(left,right)$가 SHA profile의 Merkle
compressor다. 이 직접 modulo reduction은 RFC 9380 Hash-to-Field나 \texttt{fr.Hash}와
동일하지 않으며, 256-bit digest를 Field 하나로 축약해 분포 편향이 생길 수 있다.
이 선택은 완료된 hash scaling POC를 재현하기 위한 것이며 production hash 결정으로
일반화하지 않는다. 새로운 SHA profile artifact는 생성하지 않는다.

구현 완료 전 native Go와 gnark circuit에서 같은 input tuple이 같은 Field Element를
출력하는 reference vector를 고정해야 한다. Library default를 변경하거나
gnark/gnark-crypto version을 올리면 reference vector와 8개 circuit key를 다시 생성한다.

\subsection{Azeroth BLS12-381 코드의 참고 경계}

\texttt{Azeroth Contract}의 BLS12-381 관련 Solidity 파일은 현행 zkDPP 구현체가 아니라
과거 EIP-2537 draft를 대상으로 작성된 Groth16 실험 경로다. zkDPP는 이 코드에서
다음 architecture만 참고한다.

\begin{itemize}
  \item protocol state machine과 proof verifier를 별도 contract/library 경계로 분리한다.
  \item 공통 protocol contract는 root, nullifier, tree와 registry를 관리하고,
        Relation별 verifier는 proof와 public input 검증만 담당한다.
  \item curve-specific point와 proof serialization이 verifier ABI에 반영되어야 한다.
  \item expensive group operation과 pairing은 가능한 경우 EVM precompile에 위임한다.
  \item Merkle tree는 abstract 2-to-1 node hash 경계를 두어 Go, circuit과 Contract가
        동일한 hash를 사용하게 한다.
\end{itemize}

다음 코드는 재사용하지 않는다.

\begin{itemize}
  \item \texttt{Groth16BLS12\_381.sol}: zkDPP는 Groth16이 아니라 PLONK/KZG를 사용하며,
        이 파일은 최종 EIP-2537과 다른 과거 precompile address와 calldata layout을
        사용한다.
  \item \path{BLS12_381MixerBase.sol}: 현재 \texttt{ZklayBase} constructor와 맞지 않는다.
  \item \path{Groth16BLS12_381Mixer.sol}: 현재 \texttt{\_verifyZKProof} signature와
        맞지 않으며 현 checkout의 완성된 deploy path가 아니다.
  \item \texttt{Poseidon2.sol}과 \texttt{MiMC7.sol}: 상수와 modulus가 BN254용이므로
        BLS12-381 scalar field에서 사용하는 zkDPP hash 구현으로 가져오지 않는다.
  \item Azeroth의 verifying-key storage layout, fixed public-input count,
        encrypted account, auditor encryption과 token/NFT 분기 로직을 복사하지 않는다.
\end{itemize}

zkDPP의 on-chain proof verifier는 gnark가 BLS12-381 PLONK verifying key에서 export한
Solidity verifier를 기준으로 한다. Merkle node hash용 Solidity Poseidon2 compressor는
gnark-crypto v0.20.1 BLS12-381 default parameter와 일치하도록 별도로 구현하고,
Azeroth의 hash source를 수정해 재사용하지 않는다.

Azeroth checkout에는 Groth16 setup을 실행해 key를 생성하는 code가 포함되어 있지 않다.
Deployment script는 외부에서 미리 생성된
\path{compiled/testParameter.json}의 \texttt{vk}와 \texttt{vk\_nft}를 읽어
main mixer constructor에 전달한다. 따라서 이 repository는 zkDPP의 KZG SRS 생성 방식에
대한 reference가 아니며, verifier와 protocol contract를 분리하는 architecture만 참고한다.

\subsection{Nullifier와 Recall Voucher}

\[
nf=\Hash(sk_{own},cm)
\]

\[
rv=\Hash(
\DocumentHash,q_T,\mathbf{s}_T,
addr_{Sender},addr_{Receiver},\Delta_{epoch},o_{rv}
)
\]

\[
rvnf=\Hash(rv,o_{rv})
\]

\[
\mathit{voucher}=\{
rv,\DocumentHash,q_T,\mathbf{s}_T,
addr_{Sender},addr_{Receiver},\Delta_{epoch},o_{rv}
\}
\]

$nf$는 finalized note 소비용이고, $rvnf$는 Recall Voucher resolution용이다.
$o_{rv}$는 sender와 receiver가 off-chain으로 공유한다. Hash input에는 현재 별도의
문자열 tag를 넣지 않는다. Voucher는 역할별로 다른 wallet schema를 사용하지 않는다.
Sender와 Receiver는 자기 address가 중복되더라도 위 canonical object 전체를 보관한다.
$rvnf$는 Solidity의 $\List{rvnf}$에 해당하는
\texttt{mapping(uint256 => bool)}의 key로 사용한다. 별도
\texttt{resolved[rv]} mapping은 두지 않는다.

\subsection{Public Input Ordering과 Solidity Encoding}

Protocol Specification의 각 Event에 정의된 $\mathbf{x}$와 $\mathbf{w}$는 중괄호로
표기하지만 unordered set이 아니라 ordered tuple이다.
Private witness는 가능한 경우 다음 canonical group 순서를 사용한다.

\begin{enumerate}[nosep]
  \item ownership secret과 public identifier: $sk$, $pk$
  \item note 또는 voucher 저장 순서의 $\DocumentHash$, quantity와 StateVector
  \item address, deadline과 event-specific arithmetic auxiliary
  \item Merkle path
  \item opening과 새로 뽑은 randomness
\end{enumerate}

공유하거나 중복되는 field는 한 번만 적되 각 group 내부의 object 순서를 유지한다.
Tuple에 적힌 왼쪽에서 오른쪽 순서를 gnark public witness와 Solidity verifier input
array에 동일하게 사용한다. 중첩 tuple과 고정 배열을 Field Element 단위로 펼친 최종
순서는 다음과 같다. 아래 순서는 0-based array index 순서이며 manifest의
\texttt{publicInputs}에도 그대로 기록한다.

\begin{longtable}{@{}p{0.14\textwidth}p{0.78\textwidth}@{}}
\toprule
Relation & Ordered public Field Elements \\
\midrule
\endhead
Entry & $cm_{new}$ \\
Transfer & $rt,cm_{in},nf,rv,cm_C,\Delta_{epoch}$ \\
Proceed & $rvrt,rv,rvnf,cm_{recv}$ \\
Recall & $rvrt,rv,rvnf,cm_{return}$ \\
Merge & $rt_1,rt_2,cm_1,cm_2,nf_1,nf_2,cm_{out}$ \\
Split & $rt,cm_{in},nf,cm_{out,1},cm_{out,2}$ \\
Process$^{*}$ & $\begin{aligned}
&m,n,p_{totalLoss},p_{carbonAdded},p_{recycledDelta},\\
&A_{1,total},A_{1,carbon},A_{1,recycled},
 A_{2,total},A_{2,carbon},A_{2,recycled},\\
&A_{3,total},A_{3,carbon},A_{3,recycled},\\
&rt_1,cm_{in,1},nf_1,rt_2,cm_{in,2},nf_2,\\
&rt_3,cm_{in,3},nf_3,\\
&cm_{out,1},cm_{out,2},cm_{out,3}
\end{aligned}$ \\
Exit & $rt,cm_{in},nf$ \\
\bottomrule
\end{longtable}

Process$^{*}$ 행은 기능 확인용 기본 $\ell=3,a=3$ profile의 26개 public input을 보인다.
일반 benchmark
profile에서 $a=M_{\max}=N_{\max}$이고, 공개 process-delta 개수를
$p_\ell=\ell$로 두면 public input 수는 다음과 같다.

\[
N_{public}(\ell,a)=2+p_\ell+a\ell+4a
\]

첫 2개는 $m,n$이고, 다음 $p_\ell$개는 존재하는 State coordinate의 delta다.
$\ell=3$의 delta 순서는 다음과 같다.

\[
(p_{totalLoss},p_{carbonAdded},p_{recycledDelta})
\]

마지막 값도 State coordinate와 일대일로 대응하는 public input으로 전달한다. Circuit은
recycled delta 적용식과 recycled aggregate 보존식을 함께 검사하여 유효한 proof에서
$p_{recycledDelta}=0$이 도출되게 한다. 별도의 0 assertion은 두지 않는다. 그 뒤에는
output-major $A$의 $a\ell$개 값, $(rt_i,cm_{in,i},nf_i)$ triple $a$개와 output commitment
$a$개를 차례로 놓는다.

실제 benchmark에서 사용하는 Process public input 수는 다음과 같다.

\begin{center}
\begin{tabular}{l|c|c}
\toprule
실험 & Profile & Public input 수 \\
\midrule
State scaling & $(d,a)=(32,3)$, $\ell=0,1,2,3$ & 14, 18, 22, 26 \\
Arity scaling & $(d,\ell)=(32,3)$, $a=1,2,3$ & 12, 19, 26 \\
\bottomrule
\end{tabular}
\end{center}

각 public Field Element는 canonical $0\leq x<r$ 정수 하나로 표현한다. Solidity
adapter는 이를 \texttt{uint256} 한 word로 전달하며, ABI는 각 word를 32-byte
big-endian으로 encode한다. Proof $\pi$는 public input tuple의 원소가 아니라 verifier에
별도로 전달되는 proof bytes다. Contract-only 검사값도 Protocol Specification의
$\mathbf{x}$에 포함되지 않았다면 proof public input으로 간주하지 않는다.

\section{Setup과 Artifact}

\subsection{POC Setup}

각 Relation에 대해 다음 순서를 수행한다.

\begin{enumerate}
  \item BLS12-381 scalar field와 SCS builder로 circuit을 compile한다.
  \item 해당 Relation의 \texttt{ccs}를 \texttt{unsafekzg.NewSRS(ccs)}에 전달해
        local canonical SRS와 Lagrange SRS를 생성한다.
  \item 두 SRS와 \texttt{ccs}를 PLONK Setup에 전달해 proving key와 verifying key를 생성한다.
  \item verifying key에서 Solidity \texttt{PlonkVerifier}를 export한다.
  \item CCS, 두 SRS, proving key, verifying key와 manifest를 Relation artifact
        directory에 저장하고 generated verifier는 Foundry source tree에 저장한다.
\end{enumerate}

\texttt{unsafekzg.NewSRS}는 constraint 수와 public variable 수로 필요한 SRS 크기를
계산해 canonical/Lagrange pair를 반환한다. POC에서는 각 Relation을 독립적으로 Setup한다.
이 SRS는 test용 toxic waste로 생성되므로 production security를 제공하지 않는다.

현재 setup code는 Relation마다 \texttt{NewSRS(ccs)}를 호출한다. 다만
\texttt{unsafekzg}의 in-memory cache key는 circuit identity가 아니라 curve, canonical
SRS size와 toxic value다. 따라서 같은 process 안에서 필요한 크기가 같은 Relation은
같은 SRS를 재사용하고, 크기가 다르면 새로운 random toxic value로 새 SRS를 생성한다.
Process를 다시 실행하면 memory cache도 사라지므로 새 SRS가 생성된다. 현재 artifact
directory에는 재현 편의를 위해 Relation별 SRS copy를 저장한다.

gnark v0.15.0의 required size 계산은 다음과 같다.

\[
N_{\mathrm{domain}}
=
\operatorname{NextPowerOfTwo}
\left(
\#\mathrm{constraints}+\#\mathrm{publicVariables}
\right),
\qquad
N_{\mathrm{canonical}}=N_{\mathrm{domain}}+3.
\]

Ethereum KZG Ceremony output과 현재 Ethereum blob setup은 구분해야 한다. Ceremony
specification은 서로 다른 $\tau$를 사용하는 G1 size $2^{12}$, $2^{13}$, $2^{14}$,
$2^{15}$의 네 Powers-of-Tau set을 만들었고 각 set은 G2 point 65개를 가진다. 반면
EIP-4844 blob consensus specification은 G1 monomial 4,096개, G1 Lagrange 4,096개와
G2 monomial 65개를 사용한다. 이 SRS point array는 Contract storage에 공개된 자료구조가
아니다. Ethereum client가 trusted setup을 사용하고 EVM은 point-evaluation precompile을
노출한다. 공식 근거는
\url{https://github.com/ethereum/kzg-ceremony-specs}와
\url{https://ethereum.github.io/consensus-specs/deneb/polynomial-commitments/}다.

현재 Poseidon2 Process는 domain 65,536과 canonical SRS 65,539 points가 필요하므로
최대 32,768 G1 points인 Ceremony output으로도 처리할 수 없다. Poseidon2 Entry도 gnark가
canonical 4,099 points를 요구하므로 EIP-4844의 4,096-point set을 그대로 전달할 수 없다.
큰 Ceremony set을 gnark 형식으로 변환해 Entry에 사용하는 가능성과 production SRS 선택은
별도 결정이며 현재 POC에는 적용하지 않는다.

\newpage
Artifact directory는 다음 형식을 사용한다.

\begin{verbatim}
circuits/profiles/
  state-0/
  state-1/
  state-2/
  state-3/

artifacts/poc/state-<l>/depth-<d>/<relation-variant>/
  ccs.bin
  srs-canonical.bin
  srs-lagrange.bin
  proving.key
  verifying.key
  poseidon2-vectors.json
  manifest.json

contracts/src/generated/state-<l>/depth-<d>/<relation-variant>/
  PlonkVerifier.sol
\end{verbatim}

네 \path{circuits/profiles/state-<l>} folder는 State schema와 compile constant만
정의한다. Entry부터 Exit까지의 Relation source를 네 번 복사하지 않고 공통 circuit
package를 호출한다. \texttt{relation-variant}는 Non-Process에서는 Relation의 lowercase
이름이고, Process에서는 \texttt{process-1x1}, \texttt{process-2x2} 또는
\texttt{process-3x3}이다. 새로 생성하는 profile 조합은 Benchmark Specification에 정의된
두 Poseidon2 실험에 한정한다. Setup binary는
\texttt{artifacts}에 두고 generated verifier만 Foundry가 compile하는
\texttt{contracts/src/generated}에 둔다. Generated Solidity의 contract name은 gnark
기본값인 \texttt{PlonkVerifier}를 유지한다. Foundry deployment는 동일 contract name을
파일 경로까지 포함한 fully qualified name으로 구분한다.

각 \texttt{manifest.json}은 다음 값을 갖는다.

\begin{itemize}
  \item \texttt{schemaVersion}: 정수 \texttt{1}
  \item \texttt{relationId}: Main Contract와 같은 0부터 7까지의 정수
  \item \texttt{relationName}: \texttt{entry}, \texttt{transfer}, \texttt{proceed},
        \texttt{recall}, \texttt{merge}, \texttt{split}, \texttt{process},
        \texttt{exit} 중 하나
  \item \texttt{profileId}: Non-Process는 \path{state-<l>-depth-<d>}, Process는
        \path{state-<l>-depth-<d>-process-<a>x<a>}
  \item \texttt{treeDepth}: 정수 16 또는 32, \texttt{stateLength}: 0부터 3까지의 정수
  \item Process manifest의 \texttt{processArity}: 정수 1, 2 또는 3.
        다른 Relation manifest에는 이 key를 넣지 않는다.
  \item \texttt{proofSystem}, \texttt{curve}, \texttt{gnarkVersion},
        \texttt{gnarkCryptoVersion}
  \item \texttt{constraintCount}, \texttt{publicVariableCount}
  \item \texttt{publicInputs}: Protocol Specification과 동일한 ordered name array
  \item \texttt{setupMode}: \texttt{unsafe-local}, \texttt{srsScope}: \texttt{relation}
  \item \texttt{files}: repository root 기준 상대경로와 SHA-256 checksum. 최소 key는
        \path{ccs}, \path{srsCanonical}, \path{srsLagrange},
        \path{provingKey}, \path{verifyingKey}, \path{solidityVerifier},
        \path{poseidon2Vectors}다.
  \item \texttt{sourceCommit}: 생성에 사용한 Git commit ID
\end{itemize}

Manifest의 \texttt{schemaVersion}은 artifact file format version이며 DocumentInfo의
encoding field가 아니다. Prover, verifier deployment script와 benchmark harness는
Relation ID, ordered public inputs와 checksum이 일치하지 않으면 artifact를 사용하지 않는다.
각 \texttt{files} value는 \texttt{path} string과 lowercase hex
\texttt{sha256} string을 갖는다. \texttt{sourceCommit}은 Git commit ID string이며
dirty worktree에서 생성했다면 별도 boolean \texttt{sourceDirty}를 \texttt{true}로 기록한다.

\texttt{ccs.bin}, 두 SRS와 두 key file은 해당 gnark object의 \texttt{WriteTo}가 출력한
binary bytes를 그대로 저장하고, 동일 concrete type의 \texttt{ReadFrom}으로 읽는다.
\texttt{WriteRawTo}, gzip 또는 별도 container format은 사용하지 않는다. Reader는 manifest에
고정된 gnark/gnark-crypto version과 checksum이 일치할 때만 역직렬화한다.

\texttt{poseidon2-vectors.json}은 다음 field를 갖는다.

\begin{itemize}
  \item \texttt{schemaVersion}: 정수 \texttt{1}
  \item \texttt{field}: 문자열 \texttt{BLS12-381-Fr}
  \item \texttt{compressor}: 문자열 \path{gnark-poseidon2-compress-v1}
  \item \texttt{vectors}: 각 원소가 \texttt{name}, \texttt{left}, \texttt{right},
        \texttt{output}을 갖는 array
  \item \texttt{emptyRoots}: \texttt{z16}, \texttt{z32}를 갖는 object
\end{itemize}

모든 Field Element는 \texttt{0x} prefix가 붙은 64자리 lowercase hexadecimal
canonical integer string으로 기록한다.

\subsection{BLS12-381 Solidity 재현 결과}

2026-07-16 최소 재현 실험에서 다음 조합을 확인했다.

\begin{center}
\begin{tabularx}{\textwidth}{@{}lX@{}}
\toprule
검증 항목 & 결과 \\
\midrule
gnark PLONK compile/setup/prove/verify & 성공 \\
\texttt{vk.ExportSolidity} & 성공 \\
\texttt{proof.MarshalSolidity} & 성공 \\
Foundry Prague EVM valid proof & 성공 \\
변조 public input & 거부 \\
사용 precompile & EIP-2537 G1 MSM $0x0c$, pairing $0x0f$ \\
Verifier 호출 gas & 약 329,076 gas \\
Verifier 배포 gas & 약 2,229,650 gas \\
배포 bytecode 크기 & 10,131 bytes \\
\bottomrule
\end{tabularx}
\end{center}

이 결과는 POC stack의 호환성을 확인한 것이며, 8개 실제 circuit의 gas와 proving time을
대신하지 않는다.

\section{Tree와 On-chain State}

\subsection{두 개의 Tree}

\begin{itemize}
  \item $MT$: finalized $cm$을 append하는 공급망 note tree
  \item $rvMT$: pending $rv$를 append하는 Recall Voucher tree
  \item 두 Tree는 동일한 depth $d$를 사용하는 fixed binary append-only Merkle tree다.
  \item POC와 모든 benchmark profile은 $d=32$를 사용한다.
  \item Depth scaling은 수행하지 않는다.
  \item Merkle 비용은 depth-32 commitment 한 개에 대한 membership circuit과
        Contract path/append microbenchmark로 한 번 측정한다.
\end{itemize}

\subsection{Empty Tree와 Node Hash}

두 Tree는 Azeroth \texttt{BaseMerkleTree.sol}의 logical layout을 따른다.

\[
z_0=0,
\qquad
z_{h+1}=H_{\mathrm{tree}}(z_h,z_h)
\quad h\in\{0,\ldots,d-1\}
\]

Empty root는 $z_d$다. $MT$의 leaf는 $cm$, $rvMT$의 leaf는 $rv$를 다시
leaf-hash하지 않고 그대로 사용한다.

\[
leaf_{MT}=cm\neq0,
\qquad
leaf_{rvMT}=rv\neq0
\]

$H_{\mathrm{tree}}(left,right)$는 Azeroth의 abstract \texttt{\_hash(left,right)}
경계를 유지하는 direct 2-to-1 compression function이다. Poseidon2 profile의 Native Go와
gnark circuit은 gnark-crypto/gnark가 제공하는 BLS12-381 Poseidon2 구현을 재사용한다.
완료된 historical SHA profile은 앞 절의 SHA-to-Field를 두 입력에 적용했다. Solidity에서도
선택한 profile과 같은 compressor를 주입하고 reference vector로 호환성을 검증했다.
이후 구현은 Poseidon2 compressor만 사용한다. Variable-length $H$와
$CM$이 사용하는 hasher와 Tree의 direct compressor를 혼동하지 않는다.

Bit-exact compressor는 다음 식이다.

\[
(u_0,u_1)=\mathcal{P}_{P2}(left,right),
\qquad
H_{\mathrm{tree}}(left,right)=u_1+right\pmod r
\]

두 permutation lane의 초기값은 정확히 $(left,right)$이며 별도의 zero capacity lane을
추가하지 않는다. 출력은 permutation의 오른쪽 lane에 원래 $right$를 feed-forward한
Field Element다. Native Go 구현은
\path{github.com/consensys/gnark-crypto/ecc/bls12-381/fr/poseidon2}의
\path{NewDefaultPermutation().Compress(leftBytes,rightBytes)}를 사용한다. 두 byte 입력은
각 Field Element의 canonical BLS12-381 scalar-field encoding이다. Circuit은
\path{github.com/consensys/gnark/std/permutation/poseidon2}의
\path{NewPoseidon2(api).Compress(left,right)}를 사용한다.

Pinned default permutation은 width 2, full rounds 6, partial rounds 50, degree 5이며
round-key seed string은
\texttt{Poseidon2-BLS12\_381[t=2,rF=6,rP=50,d=5]}다. Solidity compressor는 이
정의와 feed-forward를 그대로 구현해야 한다. Setup tool은 각 hash profile에 대해 최소
$H_{tree}(0,0)$, $H_{tree}(1,2)$와 $z_{32}$의 native/circuit/Solidity reference vector를 생성해
manifest와 test fixture에 저장한다.

\subsection{Append와 Contract Storage}

Leaf는 index 0부터 왼쪽에서 오른쪽으로 순차 append한다. Contract는 Azeroth처럼
실제 leaf와 append로 영향을 받은 모든 internal node를 full binary-tree heap index의
mapping에 저장한다. 기록되지 않은 storage node의 0은 해당 level의 $z_h$로 해석한다.
Valid $cm$과 $rv$는 0이 아니어야 하지만, Contract는 Azeroth와 동일하게 별도
nonzero 검사를 수행하지 않는다. Commitment 또는 voucher hash가 0이 될 확률과
Poseidon2 internal node가 0이 될 확률은 각각 약 $1/r$로 보고, 별도 existence
bitmap 없이 무시한다.

각 append는 leaf 한 개와 root까지의 $d$개 parent를 갱신한다. Tree가 $2^d$개 leaf로
가득 차면 추가 append는 revert한다.

Heap storage index는 root를 0으로 하는 0-based complete-binary-tree layout으로 고정한다.
Leaf index $j$의 storage index와 leaf에서 $h$ level 올라간 node의 storage index는 각각
다음과 같다.

\[
\operatorname{NodeIndex}(0,j)=2^d-1+j,
\qquad
\operatorname{NodeIndex}(h,j)=2^{d-h}-1+\left\lfloor\frac{j}{2^h}\right\rfloor
\]

Level $h$에서 sibling position은
$\left\lfloor j/2^h\right\rfloor\mathbin{\mathtt{xor}}1$이고 parent position은
$\left\lfloor j/2^{h+1}\right\rfloor$이다. Append와 path getter는 이 식을 함께 사용한다.

Append 함수는 삽입 전에 현재 leaf count를 zero-based $index$로 저장하고, 삽입과 root
갱신이 끝난 뒤 다음 event 중 하나를 emit한다.

\[
\mathrm{MTAppend}(cm,index,rt_{new})
\]

\[
\mathrm{RVMTAppend}(rv,index,rvrt_{new})
\]

$cm$ 또는 $rv$는 indexed event field로 두어 Client가 자신의 leaf를 검색할 수 있게
한다. Azeroth 구현은 삽입 후 증가한 \texttt{\_numLeaves}를 \texttt{index}라는 이름으로
emit하므로 실제 zero-based leaf index보다 1 크다. zkDPP는 이 동작을 그대로 복사하지
않고 삽입 전 값을 정확한 zero-based index로 emit한다.

\subsection{Membership Witness와 Root History}

Membership witness는 다음 logical object다.

\[
path=(index,sibling_0,\ldots,sibling_{d-1})
\]

$index$는 $0\leq index<2^d$인 leaf index이고, $sibling_0$은 leaf level의 sibling,
$sibling_{d-1}$는 root 바로 아래 level의 sibling이다. Circuit은
\texttt{api.ToBinary(index,d)}로 direction bit를 만들고 각 level에서
$H_{\mathrm{tree}}(left,right)$의 input ordering을 선택한다. gnark가 direction 선택
constraint를 생성하지만, Client는 index와 $d$개 sibling을 witness로 제공해야 한다.

gnark의 기본 Merkle gadget은 leaf에 FieldHasher를 한 번 적용하고 internal node에도
Merkle--Damg\aa rd hasher를 사용하므로 이 Tree와 root가 일치하지 않는다. 각 Event
circuit은 raw leaf에서 시작해 선택한 profile의 direct 2-to-1 compressor를 정확히 $d$번
적용하는 zkDPP 전용 membership verifier를 사용한다.

Contract는 $MT$와 $rvMT$의 empty root 및 append로 생성된 모든 과거 root를 각각
bool mapping에 영구 등록한다. Root-history window를 두지 않으며, spend와 resolution
proof는 해당 Tree의 accepted root 중 하나에 binding되어야 한다.

Client의 POC path 조회는 Contract view 함수로 제공한다.

\begin{verbatim}
getMTPath(uint256 index)
  returns (uint256 root, uint256[] siblings)

getRVMTPath(uint256 index)
  returns (uint256 root, uint256[] siblings)
\end{verbatim}

두 함수는 $index<leafCount$를 검사하고, 조회 시점의 current root와 leaf-to-root 순서의
sibling $d$개를 반환한다. Client는 append event에서 받은 zero-based index를 전달한다.
View 함수는 과거 root의 path를 재구성하지 않는다. 과거 root용 proof를 생성하려는
Client는 해당 시점의 path를 별도로 보관해야 한다.

Proceed와 Recall Client는 wallet이 $rv$와 함께 보관한 $index_{rv}$를
\texttt{getRVMTPath(index\_rv)}에 전달해 $rvrt$와 $path_{rv}$를 조회한다.
$index_{rv}$는 Voucher hash에 포함되는 protocol data가 아니라, Transfer 성공 시
$rvMT$ append event에서 얻어 Voucher와 연결해 보관하는 wallet metadata다.

\subsection{Logical Storage}

\[
\List{rt},\List{cm},\List{nf},
\List{rvrt},\List{rv},\List{rvnf},\Map{user},\Map{rv}
\]

POC에서 Field Element는 Solidity \texttt{uint256}으로 표현한다. 위의 logical
\texttt{List}는 순회용 array가 아니라 다음 membership mapping으로 구현한다.

\begin{itemize}
  \item $\List{cm}$, $\List{nf}$, $\List{rv}$, $\List{rvnf}$:
        \texttt{mapping(uint256 => bool)}
  \item $\Map{user}$: Solidity EOA address를 key로 하고
        \texttt{UserRegistration} struct를 value로 갖는 mapping. Struct는
        \texttt{uint256 addr}와 \texttt{uint256 pkOwn}을 갖는다. 동일 EOA의 재등록은
        기존 값을 덮어쓴다.
  \item $\List{rt}$와 $\List{rvrt}$: Tree별
        \texttt{mapping(uint256 => bool)} accepted-root set
  \item $\Map{rv}$: \texttt{mapping(uint256 => uint256)}으로 구현하는
        $rv\mapsto deadlineEpoch$
  \item 각 Tree의 internal node: \texttt{mapping(uint256 => uint256)}, leaf count:
        \texttt{uint256}
\end{itemize}

Protocol 검증에는 전체 set 순회가 필요하지 않으므로 별도 dynamic array를 두지 않는다.
Append index와 새 root는 event로 제공한다. 이 storage layout을 POC gas benchmark의
baseline으로 사용한다.

\begin{itemize}
  \item Transfer 시 Smart Contract가 current epoch를 계산하고 여기에
        $\Delta_{epoch}$을 더한 값을 $\Map{rv}[rv]$에 저장한다. 완성된 deadline
        값을 transaction input으로 받지 않는다.
  \item $EPOCH\_SIZE$는 600초다.
  \item Sender가 $\Delta_{epoch}\geq0$을 선택한다. Protocol policy에는 별도의
        maximum을 두지 않는다. Circuit은 이를 별도로 bit range check하지 않고
        voucher hash에 binding한다. Contract calldata의 $\Delta_{epoch}$은 proof public
        input과 동일한 Field Element를 \texttt{uint256}으로 받는다. Generated verifier가
        $\Delta_{epoch}<r$을 포함한 public-input canonicality를 검사하므로 Main Contract는
        이를 반복 검사하지 않는다. Proof 검증 후
        $deadlineEpoch=block.timestamp/600+\Delta_{epoch}$을 계산한다. Field modulus와
        epoch 정의에 의해 이 합은 \texttt{uint256} 범위 안에 있다.
  \item Recall은 $currentEpoch<\Map{rv}[rv]$일 때만 허용한다.
        $\Delta_{epoch}=0$이면 Recall은 처음부터 불가능하다.
  \item Recall 또는 Proceed 이후에도 $\Map{rv}[rv]$를 삭제하지 않는다.
  \item Proceed는 $rvMT$ membership으로 Transfer에서 생성된 voucher임을 검증하므로
        $\Map{rv}[rv]$ 존재 여부를 별도로 검사하지 않는다. Recall도 deadline 비교가
        존재하지 않는 mapping의 기본값 0을 자동으로 거부하므로 별도 exists 검사를 두지 않는다.
  \item Proceed와 Recall은 $rvMT$ membership proof와 별도로
        $rv\in\List{rv}$를 proof 검증 전 검사한다. 정확성을 위한 필수 membership 검사를
        대체하지 않고, 등록되지 않은 $rv$ 요청을 저렴하게 거부하는 사전 filter로 유지한다.
  \item $\List{nf}$와 $\List{rvnf}$는 중복 소비와 중복 resolution을 거부한다.
  \item accepted $MT$ root와 $rvMT$ root는 서로 다른 bool mapping으로 구현한다.
\end{itemize}

\section{Circuit 구현 요구사항}

\subsection{공통 Constraint}

필요한 Event는 다음 검사를 조합한다.

\begin{enumerate}
  \item old commitment membership과 opening binding
  \item owner address와 secret-key binding
  \item $nf$ 또는 $rvnf$ derivation
  \item quantity와 StateVector range
  \item input/output DocumentHash relation
  \item output commitment binding
  \item Event-specific conservation, proportionality 또는 process allocation
\end{enumerate}

정확한 public statement $\mathbf{x}$와 private witness $\mathbf{w}$ ordering은 Protocol
Specification의 각 Relation을 기준으로 한다. Solidity calldata ordering은 별도의
machine-readable manifest로 생성한다. 이는 새로운 Protocol Decision이 아니라
Protocol Specification의 ordering을 artifact에 전사하는 구현 작업이다.

Entry와 숫자 상태를 새로 계산하는 Transfer, Merge, Split, Process Relation은 새 output의
$q$와 StateVector coordinate를 unsigned 64-bit로 range check한다. 소비 Relation은
$MT$ 또는 $rvMT$에 append된 객체가 생성 Relation에서 이미 range를 검증받았다는
closed-state invariant를 사용하므로 input stored value의 range를 반복 검사하지 않는다.
Proceed와 Recall은 검증된 Voucher의 $q_T$와 $\mathbf{s}_T$를 변경하지 않고 새 holder용
commitment로 복원하므로 해당 range를 반복 검사하지 않는다. Exit도 새 숫자 상태를
생성하지 않으므로 input range를 반복 검사하지 않는다. Event relation에 필요한 positivity,
aggregate와 곱셈 intermediate bound는 별도로 검사한다.

\subsection{Event별 구현 범위}

\begin{longtable}{@{}p{0.14\textwidth}p{0.12\textwidth}p{0.64\textwidth}@{}}
\toprule
Event & Arity & 구현 핵심 \\
\midrule
\endhead
Entry & $0\to1$ & new commitment opening, $0<q<2^{64}$, $0\leq s_k<2^{64}$, duplicate commitment 거부, $MT$ append. 동일 DocumentHash와 fresh opening에 의한 서로 다른 commitment를 허용하며 별도 initial product semantic invariant 없음 \\
Transfer & $1\to(rv,cm)$ & old note 소비, $rv$와 change commitment 생성, deadline 등록. change를 내림 계산하고 residual은 첫 output인 $rv$에 할당 \\
Proceed & $rv\to cm$ & receiver authorization, unresolved voucher, receiver commitment 생성, $rvnf$ 기록 \\
Recall & $rv\to cm$ & sender authorization, deadline boundary, return commitment 생성, $rvnf$ 기록 \\
Merge & $2\to1$ & 동일 DocumentHash, quantity/StateVector coordinate 합, 두 input 소비 \\
Split & $1\to2$ & 동일 DocumentHash와 input 소비. 두 번째 output을 내림 계산하고 residual은 첫 번째 output에 할당 \\
Process & $a\to a$ profile & $a\in\{1,2,3\}$인 padded fixed circuit, input/output aggregate
$\mathbf{s}_{in},\mathbf{s}_{out}$, delta $P$, allocation matrix $A$, output 1 residual과
output binding \\
Exit & $1\to0$ & membership, ownership, nullifier correctness와 소비 기록 \\
\bottomrule
\end{longtable}

\subsection{Fixed-point Encoding과 Allocation Policy}

POC의 저장값 범위와 scale은 확정됐다.

\[
0\leq q<2^{64},
\qquad
0\leq s_k<2^{64},
\qquad
\mathrm{Scale}_{State}=10^9
\]

$q$는 scale 없는 integer count이고, $s_k$는 physical value에 $10^9$을 곱해
note에 저장하는 nano scale fixed-point integer다. Stored quantity와 StateVector의
64-bit range는 확정됐지만 aggregate와 cross multiplication의 intermediate bound는
Event별로 분석해야 한다.

Entry는 $0<q<2^{64}$와 $0\leq s_k<2^{64}$를 검사한다. Entry quantity는
0을 허용하지 않지만 각 $s_k=0$은 허용하며, attribute 간 별도 initial product
semantic invariant를 추가하지 않는다.

Transfer와 Split은 deterministic first-output remainder assignment를 사용한다.
각 State coordinate에서 non-first output의 quotient와 remainder를 다음 정수 관계로
검증한다.

\[
q_{floor}s_{in,k}=q_{in}s_{floor,k}+r_k,
\qquad
0\leq r_k<q_{in}
\]

첫 번째 output은 input State에서 floor output을 뺀 값으로 계산한다.

\[
s_{out,1,k}+s_{floor,k}=s_{in,k}
\]

Transfer에서는 $cm_{change}$가 floor output이고 첫 output인 $rv$가 residual을 받는다.
Split에서는 $cm_{out,2}$가 floor output이고 $cm_{out,1}$이 residual을 받는다.
Remainder는 private witness이며 public input이나 Contract storage에 추가하지 않는다.
Split은 $q_1>0$ 또는 $q_2>0$ 같은 output nonzero constraint를 추가하지 않는다.
각 output에는 공통 unsigned 64-bit range와 $q_{in}=q_1+q_2$만 적용한다.
$0\leq r_{2,k}<q_{in}$이 성립하려면 $q_{in}$은 자동으로 양수여야 하므로,
별도의 $q_{in}>0$ assert도 두지 않는다.

$q_1=0$이면 $q_2=q_{in}$이고 Relation이 $\mathbf{s}_1=\mathbf{0}$을 강제한다.
$q_2=0$이면 floor/remainder relation이 $\mathbf{s}_2=\mathbf{0}$을 강제한다.
따라서 zero-quantity output commitment에 nonzero State를 할당할 수 없다.

Field Element에는 자연스러운 대소관계가 없으므로 $0\leq r_{2,k}<q_{in}$은
64-bit bounded comparison gadget으로 구현한다. 이는 별도의 Protocol-level input
range assert를 다시 추가한다는 뜻이 아니라, field equality만으로 표현할 수 없는 정수
부등식을 bit decomposition 또는 동등한 bounded-difference constraint로 구현한다는 뜻이다.
Comparator 구현이 $q_{in}$의 bit decomposition을 내부적으로 요구하면 그 constraint는
remainder bound 검증 비용에 포함한다.

Split Client는 wallet이 보관한 $index_{in}$을 \texttt{getMTPath(index\_in)}에 전달해
$rt$와 $path_{in}$을 조회한다. $index_{in}$은 note commitment payload가 아니라 MT append
event에서 얻어 note와 함께 보관하는 wallet metadata다.

\subsection{Process}

각 Process circuit profile의 최대 arity는 다음으로 고정한다.

\[
M_{\max}=a,
\qquad
N_{\max}=a,
\qquad
a\in\{1,2,3\}
\]

기능 확인용 기본 profile은 $a=3$이다. 실제 active input/output 수는
$1\leq m\leq a$, $1\leq n\leq a$이며, 활성 slot은 index
1부터 연속된다. 별도의 public 또는 private enable vector는 사용하지 않는다. Public $m,n$에서
도출한 고정 algebraic gating expression으로 나머지 slot을 비활성화한다. 이 POC로
constraint 수, proof 생성 시간, Go verification 시간, Solidity verifier gas와 Process
transaction 전체 gas를 측정한다.

\[
E^{in}_i(m)=\mathbf{1}[i\leq m],
\qquad
E^{out}_j(n)=\mathbf{1}[j\leq n]
\]

이 값은 Client가 제공하는 witness가 아니라 circuit에 고정된 expression이다. $a=3$ profile은
다음 식을 사용할 수 있다. $2^{-1}$은 scalar field에서 2의 multiplicative inverse다.

\[
E_1(z)=1,
\qquad
E_2(z)=1-(z-2)(z-3)2^{-1},
\qquad
E_3(z)=(z-1)(z-2)2^{-1}
\]

\[
(m-1)(m-2)(m-3)=0,
\qquad
(n-1)(n-2)(n-3)=0
\]

각 input/output note의 StateVector와 aggregate를 다음처럼 구분한다.

\[
\mathbf{s}_{in}:=\sum_{i=1}^{m}\mathbf{s}_{in,i},
\qquad
\mathbf{s}_{out}:=\sum_{j=1}^{n}\mathbf{s}_{out,j}
\]

$P\in\mathbb{Z}^{\ell}$는 공정 중 발생한 coordinate별 net delta다. Profile별 State
prefix와 방향은 다음과 같다.

\[
\begin{array}{c|c|c}
\ell & P & P_{pub} \\
\hline
0 & () & () \\
1 & (-p_{totalLoss}) & (p_{totalLoss}) \\
2 & (-p_{totalLoss},+p_{carbonAdded}) & (p_{totalLoss},p_{carbonAdded}) \\
3 & (-p_{totalLoss},+p_{carbonAdded},+p_{recycledDelta}) &
(p_{totalLoss},p_{carbonAdded},p_{recycledDelta})
\end{array}
\]

$P_{pub}$의 coordinate는 public unsigned magnitude다. $\mathbf{s}_{in}$과
$\mathbf{s}_{out}$은 circuit 내부에서 계산하고 다음 relation으로 공정 변화량을
정확히 검사한다. 두 aggregate는 별도 input으로 전달하지 않는다. 기능 확인용 Relation은
$\ell=3$으로 compile한다. $\ell<3$ benchmark circuit은 StateVector prefix만 사용하고
존재하지 않는 coordinate의 delta relation을 compile하지 않는다.

Protocol Relation의 vector 표기는 다음과 같다.

\[
P:=(-p_{totalLoss},+p_{carbonAdded},+p_{recycledDelta}),
\qquad
\mathbf{s}_{out}=\mathbf{s}_{in}+P
\]

Circuit에서는 음수 Field input을 직접 받지 않고 public unsigned magnitude를 사용하므로
위 vector equality를 다음 coordinate별 equality로 구현한다.

\[
\begin{aligned}
s_{in,total}&=s_{out,total}+p_{totalLoss},\\
s_{out,carbon}&=s_{in,carbon}+p_{carbonAdded},\\
s_{out,recycled}&=s_{in,recycled}+p_{recycledDelta},\\
s_{out,recycled}&=s_{in,recycled}.
\end{aligned}
\]

마지막 두 relation을 함께 만족하면 $p_{recycledDelta}=0$이 도출된다. 별도의
$p_{recycledDelta}=0$ assertion은 추가하지 않는다.

\[
0\leq p_{totalLoss}<2^{66},
\qquad
0\leq p_{carbonAdded}<2^{66}
\]

$M_{\max}=N_{\max}=3$이고 각 active State가 64-bit이므로 aggregate는
$3(2^{64}-1)<2^{66}$이다. 이 bound는 aggregate와 public delta를 Field의 modular
wraparound 없이 정수 의미로 비교하기 위해 사용한다.

Allocation denominator는 circuit constant로 고정한다.

\[
D:=10^9
\]

$A\in\mathbb{Z}^{n\times\ell}$는 Client와 transaction이 전달하는 active matrix다.
Client와 Contract는 뒤에 $a-n$개의 zero row를 붙여 verifier용
$A^{pad}\in\mathbb{Z}^{a\times\ell}$을 구성한다. 각 active coordinate에 대해
$0\leq A^{pad}_{j,k}\leq D$와 $\sum_jA^{pad}_{j,k}=D$를 검사한다.

Verifier public input의 $A^{pad}$는 profile별 fixed $a\times\ell$ row-major 순서로
직렬화한다. 기본
$a=3,\ell=3$ profile의 순서는 다음과 같다.

\[
(A_{1,total},A_{1,carbon},A_{1,recycled},
  A_{2,total},A_{2,carbon},A_{2,recycled},
  A_{3,total},A_{3,carbon},A_{3,recycled})
\]

$j>n$인 inactive output row의 $\ell$개 값은 Client와 Contract가 모두 0으로 채운다.

\[
s_{out,k}A^{pad}_{j,k}=Ds_{out,j,k}+r_{j,k},
\qquad
0\leq r_{j,k}<D
\quad j\in\{2,\ldots,n\}
\]

\[
s_{out,1,k}+\sum_{j=2}^{n}s_{out,j,k}=s_{out,k},
\qquad
\sum_{j=1}^{n}A_{j,k}=D
\]

Process는 output 2부터 active output $n$까지 내림을 검증하고, output 1이 모든
State coordinate의 residual을 받는다. 따라서 State coordinate마다 $n-1$개의
floor relation과 remainder range check, 한 개의 aggregate conservation equality를
사용한다. POC의 $n=3$, $\ell=3$ maximum case에서는 각각 6개의 floor relation과
remainder range check, 3개의 conservation equality가 필요하다.

Inactive input의 public $rt_i,cm_i,nf_i$와 private note opening/path는 0으로 고정한다.
Inactive output의 public $cm_j$, private note opening, remainder와 $A$의 해당 row도 0으로
고정한다. Active input만 membership/nullifier/commitment relation을 검증하고 active
output만 commitment/allocation relation을 검증한다.

Main Contract의 Process ABI는 active input root, commitment, nullifier 배열, active
$n\times\ell$ allocation과 output commitment 배열을 받는다. Contract는
배열 길이에서 $m,n$을 구하고 verifier용 fixed-length public input을 trailing zero로
padding한다. $m,n$은 별도 calldata argument로 받지 않는다. Contract는 input root,
input commitment와 nullifier 배열의 길이가 같은지,
$1\leq m\leq M_{\max}$와 $1\leq n\leq N_{\max}$를 검사한다. State update에서는
active $nf$만 기록하고 active output $cm$만 $MT$에
append한다. Inactive zero slot은 calldata의 active 배열과 Tree에 포함되지 않는다.

Process Client는 $m$개의 input note/index pair와 $n$개의
$(\DocumentHash_{out,j},q_{out,j})$ output descriptor를 입력받는다. $m,n$은 이 tuple 표기의
cardinality이며 Client Scheme에서 배열 길이 계산과 arity assertion을 반복하지 않는다.
Client는 active $A\in\mathbb{F}^{n\times\ell}$만 입력받고 proof statement를 만들 때
$A^{pad}$로 zero-padding한다. Contract도 transaction의 같은 active $A$를 동일한 순서로
padding해야 하며, 두 결과가 다르면 proof verification이 실패한다.
Process는 quantity transition과 input/output
DocumentHash relation을 강제하지 않으므로 이 값들은 $P$ 또는 $A$에서 도출하지 않는다.
Relation은 각 값이 해당 output commitment에 binding되고
$1\leq q_{out,j}<2^{64}$를 만족하는지만 검증한다. Active Process output은 quantity가
0일 수 없다. 필요하지 않은 output은 active 배열과 $n$에서 제외한다.

하나의 Process transaction에서 active input $cm$은 pairwise distinct여야 한다. 이 검사는
같은 finalized note를 여러 slot에 반복해 aggregate에 중복 반영하는 것을 막기 위해 proof
verification 전에 수행한다.

Split과 Process output에는 별도 pairwise 비교를 추가하지 않는다. Proof verification 후
각 output에 대해 $\List{cm}$ 부재 확인, $\List{cm}$ 등록, $MT$ append를 index 순서로
실행한다. 앞 output이 등록된 뒤 다음 output을 검사하므로 같은 transaction 안의 중복도
기존 $\List{cm}$ 검사에서 거부된다. 중간 iteration에서 실패하면 Solidity transaction의
atomicity에 의해 앞선 mapping과 Tree update도 모두 revert된다.

POC는 on-chain recipe registry를 두지 않는다. Public $A$는 proof와 allocation을
공개적으로 binding하지만, 호출자가 선택한 $A$가 물리적으로 올바르거나 승인된
allocation rule임을 보장하지 않는다.
Public $P$도 동일하게 선언된 process delta를 proof에 binding할 뿐 실제 설비에서 발생한
변화량임을 보장하지 않는다. 이번 POC는 State accounting cost와 Relation 동작을 측정하는
범위로 제한하며, MaterialClassId, Quantity transition, Recipe commitment/root,
Recipe membership proof와 Recipe별 verifier를 구현하지 않는다.
Process input/output의 DocumentHash 사이에는 equality 또는 transition relation을 두지
않는다. 각 DocumentHash가 해당 input/output commitment에 binding되는지만 검사한다.
Process의 input/output quantity 사이에는 conservation, proportionality 또는 allocation
relation을 두지 않는다. $q_{in,i}$의 unsigned 64-bit range는 기존 note에서 계승하고
input commitment opening에 binding되는지만 검사한다. 새 $q_{out,j}$는
$1\leq q_{out,j}<2^{64}$와 output commitment binding을 검사한다. 추후 quantity transition이 필요하면
StateAllocationMatrix와 별도의 QuantityAllocationVector relation으로 추가한다.

\section{Smart Contract 구현 요구사항}

POC의 공급망 Event와 Recall resolution Event 함수는 permissionless다. Contract는
Event submitter의 \texttt{msg.sender}를 protocol $addr$ 또는 proof witness와 비교하지
않는다. \texttt{RegisterUser}는 permissionless registration이지만 호출자의
\texttt{msg.sender}를 편의용 $(addr,pkOwn)$ metadata에 연결한다. 등록 관계의
uniqueness나 hash relation은 검사하지 않고 public mapping getter만 제공한다.
Registration event는 emit하지 않는다. 따라서 registration storage gas는
측정하되, supply-chain Event별 role/administrator access-control gas는 측정하지 않는다.

\begin{itemize}
  \item 8개 Relation별 generated \texttt{PlonkVerifier}를 각각 독립 contract로 배포한다.
  \item Main zkDPP contract constructor는 8개 verifier address를 받아 고정 Relation ID에 binding한다.
  \item Main contract는 공통 \texttt{Verify(bytes,uint256[])} interface로 해당 Relation verifier를 호출한다.
  \item 함수별 실행은 on-chain precondition 검사, proof verification, state update의
        세 단계 순서를 따른다.
  \item 이미 존재하는 $cm$, $nf$, $rv$, $rvnf$와 허용되지 않은 root는 거부한다.
  \item Process의 active input commitment 중복은 proof verification 전에 거부한다.
        다중 output에는 pairwise 비교를 추가하지 않는다. Proof verification 후
        $\List{cm}$ 부재 확인, 등록과 Tree append를 output 순서대로 실행해 같은
        transaction 내부의 중복도 거부한다.
  \item proof verification이 성공하기 전 영구 state를 변경하지 않는다.
  \item proof 또는 state check 실패 시 transaction 전체가 revert되어야 한다.
  \item Merge의 두 public input commitment는 proof verification 전에
        $cm_1\neq cm_2$인지 검사한다. Public 값의 non-equality를 Circuit에서 반복 검증하지 않는다.
  \item Merge는 input별 accepted root $rt_1,rt_2$를 public input으로 받는다.
        각 $cm_i$의 membership path는 해당 $rt_i$에 binding하고, Contract는 두 root를
        각각 $\List{rt}$에서 검사한다.
  \item Recall/Proceed 이후에도 deadline history를 삭제하지 않는다.
  \item provenance graph edge event는 emit하지 않는다.
  \item $MT$와 $rvMT$ append 후에는 leaf, zero-based index와 new root를 담은
        index event를 emit한다.
\end{itemize}

Generated verifier의 ABI는 gnark v0.15.0 BLS12-381 PLONK export가 제공하는 다음
함수를 사용한다.

\[
\texttt{Verify(bytes proof, uint256[] publicInputs)}
\]

Prover는 반드시 \texttt{proof.MarshalSolidity()}의 반환 bytes를 transaction의
\texttt{proof} argument로 사용한다. Main Contract는 이 bytes를 수정하거나 다시 encode하지
않고 generated verifier에 그대로 전달한다. Public input은 위에서 고정한 ordered Field
Element를 \texttt{uint256[]}로 만든다.
Custom error의 구체 이름은 구현 단계에서 정하되 revert 조건은 이 문서의 state check를 따른다.

\subsection{Main Contract ABI}

Verifier address 배열의 Relation ID는 다음 순서로 고정한다.

\[
0=Entry,\ 1=Transfer,\ 2=Proceed,\ 3=Recall,\
4=Merge,\ 5=Split,\ 6=Process,\ 7=Exit
\]

POC Main Contract의 외부 ABI는 다음 signature를 제공한다. Parameter 이름은 설명용이며
ABI compatibility는 type과 배열 순서로 결정된다.

\begin{verbatim}
constructor(address[8] memory verifiers)

registerUser(uint256 addr, uint256 pkOwn)
entry(bytes proof, uint256 cmNew)
transfer(bytes proof, uint256 rt, uint256 cmIn, uint256 nf,
         uint256 rv, uint256 cmChange, uint256 deltaEpoch)
proceed(bytes proof, uint256 rvrt, uint256 rv, uint256 rvnf,
        uint256 cmRecv)
recall(bytes proof, uint256 rvrt, uint256 rv, uint256 rvnf,
       uint256 cmReturn)
merge(bytes proof, uint256 rt1, uint256 rt2, uint256 cm1, uint256 cm2,
      uint256 nf1, uint256 nf2, uint256 cmOut)
split(bytes proof, uint256 rt, uint256 cmIn, uint256 nf,
      uint256 cmOut1, uint256 cmOut2)
process(bytes proof, uint256[] rtIn, uint256[] processDelta,
        uint256[] allocation, uint256[] cmIn,
        uint256[] nfIn, uint256[] cmOut)
exit(bytes proof, uint256 rt, uint256 cmIn, uint256 nf)
\end{verbatim}

\texttt{process}는 배포 profile에 고정된 $\ell,a$를 사용한다. Contract는
\texttt{processDelta.length}가 $\ell$인지 검사한다. 또한
\texttt{allocation.length}가 $n\ell$인지 검사하고 뒤에 $(a-n)\ell$개의 zero를 붙여
verifier용 allocation을 구성한다. Calldata의 active array 길이로
$m,n$을 계산한 후 해당 profile의 ordered public input을 구성한다. 기본
$\ell=3,a=3$ profile에서는 26개다. Tree path getter와 public storage getter는 앞 절의
signature와 layout을 따른다.

\section{Benchmark Specification}

\subsection{비교 Profile}

benchmark depth는 $d=32$로 고정하며, 모든 축을 Cartesian product로 반복하지 않는다.
이후에는 다음 두 Poseidon2 실험을 독립적으로 수행한다.

\begin{longtable}{@{}p{0.19\textwidth}p{0.29\textwidth}p{0.44\textwidth}@{}}
\toprule
실험 & 고정값 & 변화값 \\
\midrule
\endhead
State scaling & $d=32$, Process $a=3$ & 모든 Event에서
$\ell=0,1,2,3$ \\
Process arity scaling & $d=32$, $\ell=3$ & Process만
$a=1,2,3$, $m=n=a$ \\
\bottomrule
\end{longtable}

Process 1-to-1과 2-to-2는 State 3, Depth 32에서만 생성한다. 비대칭 Process와
$(\ell<3,a<3)$의 교차 조합은 측정하지 않는다.

공통 기준점인 $(d,\ell,a)=(32,3,3)$은 한 번만 compile/setup하고 두 결과표에서
재사용한다. 중복을 제거한 artifact configuration 수는 다음과 같다.

\begin{itemize}
  \item 각 Non-Process Event: Poseidon2 State scaling 4개
  \item Process: Poseidon2 State scaling 4개와 추가 arity 1-to-1/2-to-2 2개, 총 6개
  \item 전체: $7\times4+6=34$개 benchmark configuration
  \item distinct circuit 수도 34개
\end{itemize}

여기서 configuration 수는 측정값 종류의 수가 아니다. 각 configuration마다 다음 절의
constraint, time, size와 gas 지표를 모두 기록한다.

\subsection{측정 상태}

Canonical Process profile은 input별 $rt_i$와 $P_{pub}$의 세 번째 coordinate를 포함하는
26-public-input 명세다. 이전 single-root 23-public-input implementation과 중간 24-input
명세에서 얻은 측정값은 사용하지 않는다. 최신 34개 configuration은
\path{poc-v2/benchmarks/raw.json}과 \path{poc-v2/benchmarks/summary.csv}에 기록한다.

측정은 Poseidon2, depth 32만 사용한다. Profile table의 gas는 generated verifier의
EVM \texttt{CREATE}와 \texttt{Verify} 호출을 분리해 기록한다. Storage update를 포함한
전체 Event gas는 canonical scenario의 실제 Main Contract call을
\path{poc-v2/benchmarks/canonical-event-gas.json}에 별도로 기록한다. Proof와 public-input
fixture parsing은 Event gas 측정 구간에서 제외한다.

현재 smoke run 환경은 Apple M1 Pro 10-core, 32 GB MacBook Pro, macOS 26.3이다.
시간은 단일 실행값이며 통계적 평균이 아니다. \texttt{totalAllocatedBytes}는 Go runtime의
누적 allocation 증가량이고 peak RSS가 아니다.

한 profile의 proof, proving key, verifying key와 Contract는 다른 profile과 공유하지
않는다. 측정 머신, OS, CPU, RAM, Go/gnark/Foundry 버전과 실행 옵션을 결과표에 함께
기록한다.

\subsection{측정 Layer}

\begin{longtable}{@{}p{0.21\textwidth}p{0.71\textwidth}@{}}
\toprule
Layer & 측정값 \\
\midrule
\endhead
Circuit & constraint 수, compile time, PLONK setup time, proving key 크기, verifying key 크기 \\
Prover & witness 생성 시간, proof 생성 시간, peak memory, proof 크기 \\
Off-chain verifier & Go verification 시간 \\
On-chain verifier & verifier deployment gas, \texttt{verifyProof} 단독 gas \\
Protocol transaction & 전체 gas, calldata/transaction 크기, 변경된 storage slot 수 \\
Merkle tree & append gas, path 생성 시간, sibling 수와 hash profile별 compression 수 \\
\bottomrule
\end{longtable}

Proof 생성과 Go verification은 warm-up 이후 동일 witness로 반복 측정하고 최소
median, mean, standard deviation과 p95를 기록한다. Foundry에서는 wall-clock
transaction time을 blockchain latency로 해석하지 않고 gas만 canonical on-chain
metric으로 사용한다.

\subsection{Event별 기준 Scenario}

\begin{longtable}{@{}p{0.13\textwidth}p{0.12\textwidth}p{0.67\textwidth}@{}}
\toprule
Event & 전이 & 기준 측정 Scenario \\
\midrule
\endhead
Entry & $0\to1$ & membership 없음, finalized $cm$ 한 개를 $MT$에 append \\
Transfer & $1\to2$ objects & $MT$ membership 한 개, $nf$ 한 개, $cm_{change}$를 $MT$에
append하고 $rv$를 $rvMT$에 append \\
Proceed & $rv\to1$ & $rvMT$ membership 한 개, $rvnf$ 한 개, $cm_{recv}$ 한 개 append \\
Recall & $rv\to1$ & $rvMT$ membership 한 개, deadline 검사, $rvnf$ 한 개,
$cm_{return}$ 한 개 append \\
Merge & $2\to1$ & $MT$ membership 두 개, $nf$ 두 개, output $cm$ 한 개 append \\
Split & $1\to2$ & $MT$ membership 한 개, $nf$ 한 개, output $cm$ 두 개 append.
non-divisible vector와 first-output residual 경계를 포함 \\
Process & $a\to a$ & $a\in\{1,2,3\}$별 모든 slot enabled, $MT$ membership $a$개,
$nf$ $a$개, output $cm$ $a$개 append \\
Exit & $1\to0$ & $MT$ membership 한 개와 $nf$ 한 개, Tree append 없음 \\
\bottomrule
\end{longtable}

Depth 비교는 State 3의 동일한 semantic input을 사용하고, State scaling은 앞서 정의한
prefix만 확장한다. 추가 boundary scenario는 Transfer의 zero change, $a=3$ Process에서
일부 slot만 활성화한 경우와 Recall의 deadline 경계를 포함하되 기준 비교표와 분리한다.

\section{Test Specification}

각 Event는 최소한 다음 test를 가져야 한다.

\begin{enumerate}
  \item valid witness의 Go proof 생성과 Go verification 성공
  \item 동일 proof의 Solidity verifier 성공
  \item public input 한 항목 변조 시 실패
  \item private witness relation 위반 시 proof 생성 또는 verification 실패
  \item duplicate $cm$, $nf$, $rv$, $rvnf$ 거부
  \item 잘못된 root와 expired/not-yet-valid timing boundary 거부
  \item state update 전후 tree root와 mapping 확인
  \item stored 64-bit upper bound, zero value와 Event별 intermediate bound 확인
  \item Transfer/Split의 divisible, non-divisible와 최대 residual 경계 확인
  \item Split/Process의 순차 $\List{cm}$ 등록에 의한 duplicate output 거부와
        Process duplicate input 거부
\end{enumerate}

Process는 사용 slot 수가 0, 1, maximum인 경우, non-prefix selector와 inactive slot에
nonzero witness를 넣은 경우를 추가로 검증한다. Foundry test는
\texttt{evm\_version = "prague"}를 명시한다.

Depth는 변화시키지 않는다. 별도 depth-32 membership microbenchmark는 commitment 한 개,
public root 한 개와 private index/sibling 32개를 사용하고 constraint 수, proof 생성 시간,
Go verification 시간과 Solidity verifier gas를 기록한다.

\section{Open Decisions}

다음 항목은 아직 확정되지 않았다. 구현자는 선택값을 코드에 고정하기 전에 Decision
Log를 갱신하고 이 문서를 수정해야 한다.

\begin{longtable}{@{}p{0.31\textwidth}p{0.61\textwidth}@{}}
\toprule
영역 & 미확정 질문 \\
\midrule
\endhead
Registration policy & EOA와 protocol address의 uniqueness, key rotation, overwrite 허용과
registration event 정책. POC Event 구현을 차단하지 않는 비핵심 결정 \\
Contract & custom error naming과 production deployment의 upgradeability \\
Provenance & Tree index event와 별개로 input/output commitment edge event를 도입할지 여부 \\
Production setup & canonical KZG SRS 출처, integrity 검증과 artifact 변환 \\
\bottomrule
\end{longtable}

\subsection{POC 결정과 구현의 우선순위}

Stored value와 intermediate range의 Protocol 경계는 확정됐다. $q$와 $s_k$는 unsigned
64-bit이고, Transfer/Split의 곱은 128-bit 미만이다. Process는
$s_{out,k}<3\cdot2^{64}<2^{66}$, $D=10^9<2^{30}$,
$s_{out,k}A_{j,k}<2^{96}$을 사용한다. 이를 gnark constraint와 test vector로 전사하는 작업은
추가 Protocol Decision이 아니라 implementation task다.

\begin{enumerate}
  \item 확정된 64-bit range만 적용해 Entry vertical slice를 구현한다.
  \item deterministic remainder를 적용한 Transfer/Split 식에 대해
        확정된 intermediate bound와 remainder 경계를 test vector로 검증한다.
  \item Poseidon2 Process State scaling 4개와 State-3 arity scaling 3개를 검증한다.
        공통 $(d,\ell,a)=(32,3,3)$ 결과는 중복 생성하지 않는다.
\end{enumerate}

Solidity BLS12-381 Poseidon2 compressor와 Depth-32 Merkle Tree는 사용자 정책 결정이
남은 항목은 아니지만 모든 Event가 의존하는 첫 engineering critical path다. Entry
vertical slice에서 Go, circuit과 Solidity reference vector, append root와 gas를 먼저
검증한다. Registration policy, provenance edge event와 production SRS는 POC Event 구현을
차단하지 않는다.

Poseidon2의 Field, parameter, construction, protocol input ordering, Merkle direct
compressor와 Azeroth 코드 재사용 경계는 더 이상 Open Decision이 아니다. Native Go,
gnark circuit과 Solidity 구현의 reference vector 작성은 이미 정해진 결정을 검증하는
implementation task다. \texttt{DocumentHash}의 Hash-to-Field primitive는
BLS12-381 \texttt{fr.Hash}로 확정됐다. DocumentInfo는 NFC, UTF-8와 unsigned
32-bit big-endian byte-length prefix로 고정됐고 DST는 빈 byte string이다. Public input
순서는 Protocol Specification의 $\mathbf{x}$ tuple을 따르며 Solidity에서는 canonical
Field Element를 \texttt{uint256} word로 전달한다.

\section{추후 고려 사항}

이 절에는 사용자가 직접 추후 고려 대상으로 지정한 항목만 기록한다. Open Decisions의
항목을 구현 우선순위나 작성자의 판단으로 이 절에 이동하지 않는다.

\subsection{Policy/Recipe 기반 Process}

사후 감사 의존을 줄이려면 Process가 input/output Material class, Quantity transition,
State delta와 State allocation policy를 Recipe로 검증해야 한다. Recipe별 circuit은
구현이 직접적이지만 verifier 선택으로 공정을 노출한다. 하나의 공통 circuit과 hidden
Recipe membership proof는 Recipe 승인 주체, registry governance, data availability와
path 획득 문제를 추가한다. 이 trust/privacy model이 확정되기 전에는 POC에 도입하지
않는다.

현재 구현은 public $P,A$를 사용하는 generic State Allocation Process를 유지한다.
이 구현은 선언된 transition의 consistency와 공개 가능성을 검증하지만, $P,A$의 물리적
정당성이나 Material/Quantity process semantics를 보장하지 않는다.

\section{변경 규칙}

\begin{enumerate}
  \item Protocol 의미가 바뀌면 Protocol Specification을 먼저 변경한다.
  \item 구현 선택이 바뀌면 Development Specification과 test acceptance를 변경한다.
  \item 선택 이유, 반려안과 실험 근거는 Decision Log에 기록한다.
  \item Paper Scheme과 HTML은 위 문서에서 파생하며 source of truth로 사용하지 않는다.
\end{enumerate}

\end{document}
